\section*{Capítulo \ref{ch:g3:food-chains} --- Cadeias alimentares}

\begin{solution}{exo:g3:food-chains:1}
``\omterm{def:g3:food-chains:chain}{É comido por}''. Ela aponta do comido para o comedor --- seguindo o
alimento.
\end{solution}

\begin{solution}{exo:g3:food-chains:2}
Uma \omterm{def:g1:plants-around-us:plant}{planta} (ou outro \omterm{def:g1:living-or-not:living}{ser vivo} que fabrica o próprio alimento).
Porque só as \omterm{def:g1:plants-around-us:plant}{plantas} se alimentam sem comer ninguém --- a partir da
luz, da água e do ar; todo elo \omterm{def:g1:animals-around-us:animal}{animal} vive desse primeiro elo verde,
diretamente ou por meio de outros.
\end{solution}

\begin{solution}{exo:g3:food-chains:3}
Capim $\rightarrow$ gafanhoto $\rightarrow$ lagarto $\rightarrow$
gavião.
\end{solution}

\begin{solution}{exo:g3:food-chains:4}
As setas estão ao contrário: capim $\rightarrow$ coelho
$\rightarrow$ raposa --- o capim é comido pelo coelho, que é comido
pela raposa.
\end{solution}

\begin{solution}{exo:g3:food-chains:5}
\omterm{def:g1:plants-around-us:parts}{Folhas} de carvalho $\rightarrow$ \omterm{ex:g2:animals-grow-up:butterfly}{lagarta} $\rightarrow$ chapim
$\rightarrow$ gavião-miúdo.
\end{solution}

\begin{solution}{exo:g3:food-chains:6}
A alface é comida pelo caracol; o ouriço come o caracol.
\end{solution}

\begin{solution}{exo:g3:food-chains:7}
Por exemplo: \omterm{def:g1:plants-around-us:plant}{plantas} aquáticas $\rightarrow$ girino $\rightarrow$
peixinho $\rightarrow$ garça.
\end{solution}

\begin{solution}{exo:g3:food-chains:8}
Porque o alimento do alimento do alimento dele é o capim: menos
gafanhotos no capim marrom quer dizer menos lagartos, e menos
lagartos quer dizer caça magra para o gavião. Uma mudança no
primeiro elo viaja por toda a cadeia acima.
\end{solution}

\begin{solution}{exo:g3:food-chains:9}
Envenenar os caracóis mata de fome quem os come: o ouriço vai embora
ou morre. Mais tarde, com o ouriço --- o principal inimigo dos
caracóis --- fora do caminho, as lesmas e os caracóis que sobraram
se multiplicam com mais força do que antes: tirar elos viaja tanto
para cima da cadeia (ouriço com fome) quanto de volta para baixo
(caracóis sem caçador).
\end{solution}

\begin{solution}{exo:g3:food-chains:10}
Verdadeiro --- passando pelo meio da cadeia. O lagarto come
gafanhotos, e os gafanhotos comem capim: se o capim falha, a falha
viaja para cima, elo por elo, do capim ao gafanhoto e ao lagarto. O
lagarto nunca encosta no capim, e mesmo assim depende dele.
\end{solution}
